\chapter{Bytes and Bytearray}

\textbf{What this chapter gives you:} the two types Python uses for raw binary data. \code{bytes} is
read-only and \code{bytearray} can be edited in place. You will meet them every time you send an image to
an API, read a streamed response, hash a prompt, or look at what a tokenizer really sees.

%RELATED: 4, 8, 17, 18

% ------------------------------------------------------------
\section{Bytes}

A \code{bytes} object is a fixed row of small integers, each from 0 to 255. Think of beads on a wire:
all the same size, no meaning attached. A \code{str} is different: it is a row of characters.
Python never mixes the two on its own. You convert on purpose.

\begin{keyidea}[: a byte is just a number]
A byte holds one integer in the range 0--255. Python \emph{prints} some bytes as letters (\code{b"A"})
because that is easier to read. The data is still the number 65.
\end{keyidea}

%MD<
% ```text
% str    :   c     a     f     é          4 characters
%            |     |     |     |
% encode     v     v     v     v  v        .encode("utf-8")
% bytes  :  63    61    66    c3 a9        5 bytes (shown in hex)
%            ^                               .decode("utf-8")
% ```
%MD>
\begin{center}
\begin{tikzpicture}[ch/.style={draw=cgold,fill=cgold!12,minimum width=10mm,minimum height=7mm,font=\ttfamily\small},
  by/.style={draw=cblue,fill=cblue!8,minimum width=10mm,minimum height=7mm,font=\ttfamily\small},
  lab/.style={font=\sffamily\scriptsize,text=cgray},>=Stealth]
  \node[ch] (c1) at (0,1.6) {c};
  \node[ch] (c2) at (1,1.6) {a};
  \node[ch] (c3) at (2,1.6) {f};
  \node[ch,minimum width=20mm] (c4) at (3.5,1.6) {\'e};
  \node[by] (b1) at (0,0) {63};
  \node[by] (b2) at (1,0) {61};
  \node[by] (b3) at (2,0) {66};
  \node[by] (b4) at (3,0) {c3};
  \node[by] (b5) at (4,0) {a9};
  \foreach \i in {1,2,3} \draw[->,cgray] (c\i) -- (b\i);
  \draw[->,cgray] (c4.south west) -- (b4.north);
  \draw[->,cgray] (c4.south east) -- (b5.north);
  \node[lab,anchor=east] at (-0.8,1.6) {\code{str}, 4 chars};
  \node[lab,anchor=east] at (-0.8,0) {\code{bytes}, 5 bytes};
  \node[lab,anchor=west] at (5.0,1.6) {text: what people read};
  \node[lab,anchor=west] at (5.0,0.8) {\code{.encode("utf-8")} down, \code{.decode("utf-8")} up};
  \node[lab,anchor=west] at (5.0,0) {numbers: what is stored and sent};
\end{tikzpicture}
\end{center}

One character can take more than one byte. That is why \code{len(text)} and \code{len(data)} can differ.

\subsection{Create with quotes}

Put a \code{b} in front of the quote. Only plain ASCII characters are allowed inside. For any other
value, write two hex digits after a backslash-x.

\begin{pycode}
>>> raw = b"GET\x00\x01\xff"
>>> raw
b'GET\x00\x01\xff'
>>> type(raw)
<class 'bytes'>
>>> len(raw)
6
>>> b"\x41\x42\x43"
b'ABC'
\end{pycode}

\subsection{Create with \texttt{bytes()}}

Give \code{bytes()} an iterable of integers (a list, a tuple, a \code{range}), or a single integer.

\begin{pycode}
>>> bytes([72, 105, 33])
b'Hi!'
>>> bytes(3)
b'\x00\x00\x00'
>>> len(bytes(range(256)))
256
>>> bytes([256])
Traceback (most recent call last):
  ...
ValueError: bytes must be in range(0, 256)
\end{pycode}

\begin{trap}[: \texttt{bytes(3)} is not \texttt{b"3"}]
A single integer means \emph{length}, not value. \code{bytes(3)} gives three zero bytes. To get one byte
whose value is 3, pass a list: \code{bytes([3])}.
\end{trap}

\subsection{Create from a hex string}

\code{bytes.fromhex()} reads pairs of hex digits and ignores spaces between the pairs. A lone digit is an
error. It is a \emph{class method}: you call it on \code{bytes} itself, not on a value (Chapter 11).
\code{.hex()} goes the other way.

\begin{pycode}
>>> bytes.fromhex("48 69 21")
b'Hi!'
>>> bytes.fromhex("4 8")
Traceback (most recent call last):
  ...
ValueError: non-hexadecimal number found in fromhex() arg at position 1
>>> b"Hi!".hex()
'486921'
\end{pycode}

\subsection{Encode and decode}

Text becomes bytes with \code{.encode()}. Bytes become text with \code{.decode()}. You must name the
encoding; use \code{"utf-8"} unless you have a reason not to. Chapter 17 covers the details.

\begin{pycode}
>>> text = "caf\xe9"
>>> data = text.encode("utf-8")
>>> data
b'caf\xc3\xa9'
>>> len(text), len(data)
(4, 5)
>>> data.decode("utf-8") == text
True
>>> data.decode("ascii")
Traceback (most recent call last):
  ...
UnicodeDecodeError: 'ascii' codec can't decode byte 0xc3 in position 3: ...
\end{pycode}

\begin{gamerule}
Always pass the encoding by name and decode with the same one you encoded with. Text on the wire (JSON,
HTTP bodies, files) is almost always UTF-8.
\end{gamerule}

\subsection{Index, slice, combine, repeat}

Bytes are a sequence, so \code{[ ]}, slices, \code{+}, \code{*}, \code{in}, and \code{for} all work.
The catch is what you get back.

\begin{typecheck}[: indexing vs slicing]
\begin{center}
\begin{tabularx}{\linewidth}{@{}l l l X@{}}
\toprule
\textbf{Expression} & \textbf{Result} & \textbf{Type} & \textbf{Note} \\
\midrule
\code{b[2]}       & \code{97}      & \code{int}   & one position gives a \emph{number} \\
\code{b[2:3]}     & \code{b"a"}    & \code{bytes} & a slice always gives \code{bytes} \\
\code{b + b}      & joined         & \code{bytes} & new object, both sides must be \code{bytes} \\
\code{b * 2}      & repeated       & \code{bytes} & new object \\
\code{for x in b} & \code{x}       & \code{int}   & the loop walks numbers \\
\code{97 in b}    & \code{True}    & \code{bool}  & an int tests for one byte \\
\code{b"a" in b}  & \code{True}    & \code{bool}  & bytes test for a run of bytes \\
\bottomrule
\end{tabularx}
\end{center}
\end{typecheck}

\begin{pycode}
>>> b = b"\xff\x01abcdef"
>>> b[2]
97
>>> chr(b[2])
'a'
>>> b[2:3]
b'a'
>>> b[-3:]
b'def'
>>> b[0:2] + b"!!"
b'\xff\x01!!'
>>> b"ab" * 2
b'abab'
>>> list(b"Hi")
[72, 105]
>>> b"abc" == "abc"
False
>>> b"abc" + "def"
Traceback (most recent call last):
  ...
TypeError: can't concat str to bytes
\end{pycode}

A useful real-world check is the first few bytes of a file, its ``magic number'':

\begin{pycode}
>>> data = b"%PDF-1.7 ...rest of file..."
>>> data[:4] == b"%PDF"
True
>>> data.startswith(b"%PDF")
True
\end{pycode}

\subsection{Bytes are immutable}

You cannot assign to a position. \code{b1 += b2} looks like an edit, but Python builds a \emph{new} bytes
object and moves the name \code{b1} onto it. The old one is dropped.

\begin{pycode}
>>> data = bytes([1, 2, 3])
>>> data[0] = 9
Traceback (most recent call last):
  ...
TypeError: 'bytes' object does not support item assignment
>>> before = data
>>> data += bytes([4])
>>> data
b'\x01\x02\x03\x04'
>>> before
b'\x01\x02\x03'
>>> data is before
False
\end{pycode}

% ------------------------------------------------------------
\section{Bytearray}

\code{bytearray} is the editable twin of \code{bytes}. Every read operation works the same. On top of
that you can change values, change the length, insert, and delete. Use it when you build binary data
piece by piece.

%MD<
% ```text
% bytes      x = b"ab"          x += b"!"
%   before:  x --> [61 62]      x --> [61 62 21]   (a NEW object)
%                  [61 62] is dropped
%
% bytearray  x = bytearray(b"ab")     x.extend(b"!")
%   before:  x --> [61 62]            x --> [61 62 21]   (the SAME object grew)
% ```
%MD>
\begin{center}
\begin{tikzpicture}[box/.style={draw=cblue,fill=cblue!8,minimum height=7mm,font=\ttfamily\small,inner sep=3pt},
  old/.style={draw=cgray,dashed,minimum height=7mm,font=\ttfamily\small,inner sep=3pt,text=cgray},
  nm/.style={draw=cgold,fill=cgold!15,rounded corners=2pt,font=\ttfamily\small,inner sep=3pt},
  lab/.style={font=\sffamily\scriptsize,text=cgray},>=Stealth]
  % bytes panel
  \node[lab,anchor=west] at (-0.4,2.5) {\textbf{bytes}: \code{x += b"!"}};
  \node[nm] (x1) at (0,1.2) {x};
  \node[old] (o1) at (2.2,1.9) {61 62};
  \node[box] (n1) at (2.6,0.5) {61 62 21};
  \draw[->,cgray,dashed] (x1) -- node[lab,above,sloped]{before} (o1);
  \draw[->,thick,cblue] (x1) -- node[lab,below,sloped]{after} (n1);
  \node[lab,anchor=west] at (3.4,1.9) {dropped};
  \node[lab,anchor=west] at (4.1,0.5) {new object};
  % bytearray panel
  \node[lab,anchor=west] at (7.6,2.5) {\textbf{bytearray}: \code{x.extend(b"!")}};
  \node[nm] (x2) at (7.8,1.2) {x};
  \node[box] (n2) at (10.4,1.2) {61 62 21};
  \draw[->,thick,cblue] (x2) -- (n2);
  \node[lab,anchor=west] at (9.2,0.5) {same object, grown in place};
\end{tikzpicture}
\end{center}

\subsection{Create with \texttt{bytearray()}}

\begin{pycode}
>>> bytearray()
bytearray(b'')
>>> bytearray([1, 2, 3, 255])
bytearray(b'\x01\x02\x03\xff')
>>> bytearray(4)
bytearray(b'\x00\x00\x00\x00')
>>> bytearray("abc", "ascii")
bytearray(b'abc')
>>> bytearray(b"from bytes")
bytearray(b'from bytes')
\end{pycode}

With a string you must give the encoding as the second argument. \code{bytearray("abc")} alone is a
\code{TypeError}.

\subsection{Read: index and slice}

Same rules as \code{bytes}, with one difference. A slice of a \code{bytearray} is a \code{bytearray}.

\begin{pycode}
>>> ba = bytearray(b"\xff\x01abcdef")
>>> ba[2]
97
>>> ba[4:6]
bytearray(b'cd')
\end{pycode}

\subsection{Change one byte or many}

\begin{pycode}
>>> ba = bytearray([1, 2, 3, 255])
>>> ba[1] = 127
>>> ba
bytearray(b'\x01\x7f\x03\xff')
>>> ba = bytearray(b"\xff\x01abcdef")
>>> ba[2:5] = b"oops"
>>> ba
bytearray(b'\xff\x01oopsdef')
>>> ba[1] = 300
Traceback (most recent call last):
  ...
ValueError: byte must be in range(0, 256)
\end{pycode}

The slice assignment above replaced three bytes with four, so the array grew. Strings cannot do this.

\subsection{\texttt{insert()}, \texttt{append()}, \texttt{extend()}}

\begin{pycode}
>>> ba = bytearray(b"abc")
>>> ba.insert(0, 2)
>>> ba.append(4)
>>> ba.extend(b"!!!")
>>> ba
bytearray(b'\x02abc\x04!!!')
>>> ba.append(b"\x04")
Traceback (most recent call last):
  ...
TypeError: 'bytes' object cannot be interpreted as an integer
\end{pycode}

\begin{typecheck}[: what each method wants]
\begin{center}
\begin{tabularx}{\linewidth}{@{}l l l X@{}}
\toprule
\textbf{Method} & \textbf{Argument} & \textbf{Returns} & \textbf{Effect} \\
\midrule
\code{ba.append(n)}      & one \code{int} 0--255  & \code{None} & adds one byte at the end \\
\code{ba.extend(bs)}     & bytes-like or ints     & \code{None} & adds many bytes at the end \\
\code{ba.insert(i, n)}   & index, \code{int}      & \code{None} & adds one byte before position \code{i} \\
\code{ba.pop()}          & optional index         & \code{int}  & removes and returns one byte \\
\code{ba.replace(a, b)}  & two bytes-like         & \code{bytearray} & a \emph{new} array; \code{ba} is unchanged \\
\bottomrule
\end{tabularx}
\end{center}
\end{typecheck}

\begin{trap}[: \texttt{replace()} does not edit in place]
On a \code{bytearray}, \code{.replace()} returns a changed \emph{copy}, like \code{str.replace()}.
If you ignore the return value, nothing happens.
\end{trap}

\begin{pycode}
>>> ba = bytearray(b"\xff\x01abcdef")
>>> new = ba.replace(b"\x01ab", b"XYZ")
>>> new
bytearray(b'\xffXYZcdef')
>>> ba
bytearray(b'\xff\x01abcdef')
\end{pycode}

\subsection{Combine, repeat, freeze}

\begin{pycode}
>>> bytearray(b"\x01\x02") + bytearray(b"\x03\x04")
bytearray(b'\x01\x02\x03\x04')
>>> bytearray(b"\x01\x02") * 2
bytearray(b'\x01\x02\x01\x02')
>>> frozen = bytes(bytearray(b"hi"))
>>> type(frozen)
<class 'bytes'>
>>> hash(bytearray(b"hi"))
Traceback (most recent call last):
  ...
TypeError: unhashable type: 'bytearray'
\end{pycode}

\code{bytes(some\_bytearray)} makes an immutable copy. Because a \code{bytearray} can change, it is
unhashable, so it cannot be a dictionary key or a set member (Chapter 9). A \code{bytes} can.

\begin{typecheck}[: bytes vs bytearray vs str]
\begin{center}
\begin{tabularx}{\linewidth}{@{}l c c c X@{}}
\toprule
 & \textbf{mutable} & \textbf{hashable} & \textbf{\code{x[0]} gives} & \textbf{slice gives} \\
\midrule
\code{bytes}     & \xmark & \cmark & \code{int}  & \code{bytes} \\
\code{bytearray} & \cmark & \xmark & \code{int}  & \code{bytearray} \\
\code{str}       & \xmark & \cmark & \code{str}  & \code{str} \\
\bottomrule
\end{tabularx}
\end{center}
\end{typecheck}

% ------------------------------------------------------------
\section{Where bytes show up in AI work}

You will not write much byte-level code, but the types appear at every boundary where data leaves
your program.

\begin{center}
\begin{tabularx}{\linewidth}{@{}l X@{}}
\toprule
\textbf{Situation} & \textbf{What you hold} \\
\midrule
Reading an image or PDF to send to a vision model & \code{bytes} from \code{open(path, "rb")} \\
Streaming reply arrives in pieces                  & \code{bytes} chunks, collected in a \code{bytearray} \\
Sending an image as JSON                           & \code{base64} turns \code{bytes} into \code{bytes} of ASCII, then \code{.decode()} to \code{str} \\
Cache key for a prompt                             & \code{hashlib} takes \code{bytes}, so \code{.encode()} first \\
Byte-level tokenizers                              & the model sees UTF-8 bytes, so one emoji can be several tokens \\
\bottomrule
\end{tabularx}
\end{center}

Two short examples, both from the standard library:

\begin{pycode}
>>> import base64, hashlib
>>> base64.b64encode(b"Hi!")
b'SGkh'
>>> base64.b64encode(b"Hi!").decode("ascii")
'SGkh'
>>> hashlib.sha256(b"prompt").hexdigest()[:8]
'cf07194e'
\end{pycode}

Streamed replies can split a multi-byte character across two chunks. Collect the raw bytes first and
decode once at the end:

\begin{pycode}
chunks = [b'{"name": "caf\xc3', b'\xa9", "tokens": 4}']
buf = bytearray()
for chunk in chunks:
    buf.extend(chunk)
print(len(buf), "bytes")
print(buf.decode("utf-8") == '{"name": "caf\xe9", "tokens": 4}')
\end{pycode}
\begin{rawcode}
30 bytes
True
\end{rawcode}

Decoding each chunk alone would fail, because the two halves of \code{c3 a9} are in different chunks.

\begin{boss}[: bytes and str look alike but are not]
The same-looking value behaves differently. \code{b"abc"[0]} is the number 97. \code{"abc"[0]} is the
string \code{"a"}. \code{b"abc" == "abc"} is \code{False}, and \code{+} between them is a
\code{TypeError}.\\[2pt]
\textbf{Strategy:} ask ``is this data text or raw bytes?'' at every boundary. Decode as soon as bytes
come in and encode as late as possible before they leave. Inside your program, work with \code{str}.
\end{boss}

\begin{practice}
\begin{enumerate}
\item Build a \code{bytes} value holding the numbers 1 to 5. What does \code{len()} return?
\item What is the difference between \code{bytes(5)} and \code{bytes([5])}?
\item Given \code{b = b"hey"}, what are the type and value of \code{b[0]} and of \code{b[0:1]}?
\item Make a \code{bytearray} from the numbers 1 to 5, then append 6 and set the first byte to 0.
\item Why does \code{ba.replace(b"a", b"b")} seem to do nothing when you ignore its result?
\item A friend encodes a string as UTF-8 and decodes it as ASCII and gets an error. Why?
\end{enumerate}
\end{practice}

\begin{achievement}{Byte Wrangler}
You can now:
\begin{itemize}
\item create \code{bytes} and \code{bytearray} from literals, lists, hex strings, and text;
\item predict whether an index or slice returns an \code{int} or a sequence;
\item convert between \code{str} and \code{bytes} with \code{.encode()} and \code{.decode()};
\item pick \code{bytearray} when data must grow or change, and freeze it with \code{bytes()};
\item recognise where bytes appear in API, file, and tokenizer work.
\end{itemize}
\end{achievement}
